\chapter{Background}
\label{background}

\section{Notation}
\subsection{Mathematical typography}
Presenting mathematics in a text is quite technical and requires lots of \LaTeX\; code. The first thing you must be aware of is that all mathematics must be placed in math mode to get the correct font. There is a huge difference between $n$ (mathematical symbol) and n (letter). Letters in math mode will always be in math font, so you must specify text: $x=2 and x=3$ is a typical error, but either split it in two math modes $x=2$ and $x=3$ or add text inside $x=2 \text{ and } x=3$. 

Please note that mathematics should be included as part of a sentence. It could be inline, like $x=2$, but then you must avoid line breaks inside the expression. If a mathematical expression is too tall or long for a single line, or it should stand out, it should be displayed centered on a separate line. For instance we have, using the equation environment, that
\begin{equation*}
    \iint_D \dd x \diff y
    =
    \int_0^{2\pi} \int_0^t \rho \diff \rho \diff t
    =
    \frac{4}{3} \pi^3,
    \label{eq:integration}
\end{equation*}
which could also be written in display style as 
\[\iint_D \dd x \diff y
    =
    \int_0^{2\pi} \int_0^t \rho \diff \rho \diff t
    =
    \frac{4}{3} \pi^3.\]
Notice the differential in \cref{eq:integration} typed with the command \texttt{\textbackslash diff}. You can cleverly refer to your equation with \texttt{\textbackslash cref}: \cref{{eq:integration}} or \texttt{\textbackslash vref}: \vref{{eq:integration}} to get the page number if it is on a different page.

\subsection{Using symbols}
 There are two golden rules that applies to the use of symbols in text.
\begin{enumerate}
    \item Don’t introduce unnecessary symbols
    \item Define all symbols in the text before they are used
\end{enumerate}

There are three typical mistakes that are made all the time - even experienced researchers do this, and it could quickly escalate into a nightmare if you don't notice it early in the process.
  \begin{itemize}
      \item One thing is referred to by two different symbols
        \item One symbol means two things
    \item Breaking the golden rules
  \end{itemize}

When you choose notation, please follow standards in your field.



\section{Theory}
We now state a theorem from {\cite{AM69}} that has a name and will be used later on. In the code, notice the braces around the precise citation inside the theorem environment.
\begin{theorem}[Dedekind's theorem {\cite[95]{AM69}}]
    \label{thm:dedekind}
    Let \( A \) be a Noetherian domain of dimension one. Then the following are equivalent:
    \begin{enumerate}[leftmargin=30pt] %%Bug: something is wrong with enumerate inside theorem-environment.
        \item \( A \) is integrally closed;
        \item Every primary ideal in \( A \) is a prime power;
        \item Every local ring \( A_\mathfrak{p} \) \( (\mathfrak{p} \neq 0) \) is a discrete valuation ring.
    \end{enumerate}
\end{theorem}
\noindent Notice that the text is in italics in the theorem environment. Sometimes it is convenient to use the command \texttt{\textbackslash noindent} to remove the paragraph indentation immediately after an environment.
\begin{proof}
This is what a proof environment looks like. Notice the square $\square$ that marks the end of the proof. You may change this symbol, but stick to standards, e.g., \enquote{QED} or $\blacksquare$.
\end{proof}


We could have talked about \emph{Fermat's last theorem} and cited \cite[Theorem~1.1]{Wil93}. Then we want to cite \cite{Vie08}, which shows that different values of \texttt{sorting} in numeric styles in \texttt{biblatex} gives citations numbers in order of appearance or alphabetic.
